\documentclass[FM,DP]{tulthesis}

\usepackage{polyglossia}
\setdefaultlanguage{czech}

\usepackage{makeidx}
\makeindex

%\usepackage{xunicode}
%\usepackage{xltxtra}

% už jsou v tul.sty
%\usepackage{fancyhdr}
%\usepackage{graphicx}
%\usepackage{ifthen}
%\usepackage{titlesec}
%\usepackage{xcolor}
\usepackage{tul}

\TULtitle{Implementace vybraných algoritmů zpracování obrazu v hradlovém poli}{Implementation of selected image processing algorithms in a gate array}

\TULprogramme{N0613}{Informační technologie}{}

\TULbranch{A140028}{Výpočetní systémy}{}
\TULauthor{Bc. Karel Najman}

\TULsupervisor{Ing. Martin Rozkovec, Ph.D.}

\TULyear{2025/2026}

% Použití bibLateXu, pracuje s ISO stylem
% BibLaTeX settings, works with ISO style
\usepackage[ 
    backend=biber
    ,style=iso-numeric
    %,style=numeric
    %,sortlocale=cs_CZ
    ,sorting=nyt %seřazeno jméno, rok, titul
    ,autolang=other
    ,bibencoding=UTF8
    %,urldate=edtf
    ,maxcitenames=2 %maximum v textu citovaných jmen
    ,maxbibnames=3 %maximum v seznamu vyjmenovaných autorů
    ]{biblatex}
\addbibresource{citace.bib}% vložení seznamu literárních zdrojů v bib formátu

% Úprava iso-numeric.bbx v souladu s požadavky TUL hranaté závorky v číslovaném seznamu / Modification of iso-numeric.bbx in accordance with TUL requirements of square brackets in a numbered list
\DeclareFieldFormat{labelnumberwidth}{\mkbibbrackets{#1}}

\usepackage{csquotes} %užití biblatexu hlasí warnings, důvodem může být použití českých uvozovek v citacích!
\urlstyle{same} %sazba url odkazů stejným fontem jako ostatní text, řešení problémů v zalamování hypertextových odkazů v citacích

\usepackage{listings} %package for sourcode typesetting
%řešení problémů češtiny např. při sazbě zdrojových kódů / czech special characters in sourcecode typesetting
\makeatletter
\lst@InputCatcodes
\def\lst@DefEC{%
 \lst@CCECUse \lst@ProcessLetter
  ^^80^^81^^82^^83^^84^^85^^86^^87^^88^^89^^8a^^8b^^8c^^8d^^8e^^8f%
  ^^90^^91^^92^^93^^94^^95^^96^^97^^98^^99^^9a^^9b^^9c^^9d^^9e^^9f%
  ^^a0^^a1^^a2^^a3^^a4^^a5^^a6^^a7^^a8^^a9^^aa^^ab^^ac^^ad^^ae^^af%
  ^^b0^^b1^^b2^^b3^^b4^^b5^^b6^^b7^^b8^^b9^^ba^^bb^^bc^^bd^^be^^bf%
  ^^c0^^c1^^c2^^c3^^c4^^c5^^c6^^c7^^c8^^c9^^ca^^cb^^cc^^cd^^ce^^cf%
  ^^d0^^d1^^d2^^d3^^d4^^d5^^d6^^d7^^d8^^d9^^da^^db^^dc^^dd^^de^^df%
  ^^e0^^e1^^e2^^e3^^e4^^e5^^e6^^e7^^e8^^e9^^ea^^eb^^ec^^ed^^ee^^ef%
  ^^f0^^f1^^f2^^f3^^f4^^f5^^f6^^f7^^f8^^f9^^fa^^fb^^fc^^fd^^fe^^ff%
% české znaky
  ^^^^010c^^^^010d^^^^010e^^^^010f^^^^011a^^^^011b^^^^0147^^^^0148%
  ^^^^0158^^^^0159^^^^0160^^^^0161^^^^0164^^^^0165%
  ^^^^016e^^^^016f^^^^017d^^^^017e%
  ^^00}
\lst@RestoreCatcodes
\makeatother

%moje balíčky
\usepackage{amsmath}   % Pro matematiku
\usepackage{booktabs}  % Pro tabulky
\usepackage{listings}
\usepackage{tikz}      % Pro vektorové nákresy a diagramy
\usepackage{caption}   % Pro popisky plovoucích prostředí
\usepackage[printonlyused, withpage]{acronym} % Pro seznam zkratek

\begin{document}

%\ThesisStart{male}
\ThesisStart{DP_zadani_a_prohlaseni_KN.pdf}

\begin{abstractCZ}
Český abstrakt
\end{abstractCZ}

\begin{keywordsCZ}
Klíčová slova česky
\end{keywordsCZ}

\vspace{2cm}

\begin{abstractEN}
English abstract
\end{abstractEN}

\begin{keywordsEN}
keywords in English
\end{keywordsEN}

\begin{acknowledgement}
\end{acknowledgement}

\tableofcontents

\clearpage

%\begin{abbrList}
\chapter*{Seznam zkratek}
\begin{acronym}[MPSoC] % To v závorce určuje odsazení (dej sem nejdelší zkratku)
    \acro{AMD}{Advanced Micro Devices}
    \acro{APU}{Application Processing Unit}
    \acro{AXI}{Advanced eXtensible Interface}
    \acro{BRAM}{Block Random Access Memory}
    \acro{DMA}{Direct Memory Access}
    \acro{DSP}{Digital Signal Processing}
    \acro{FPGA}{Field Programmable Gate Array}
    \acro{HLS}{High-Level Synthesis}
    \acro{IP}{Intellectual Property}
    \acro{LUT}{Look-Up Table}
    \acro{MIPI}{Mobile Industry Processor Interface}
    \acro{MPSoC}{Multiprocessor System on Chip}
    \acro{PL}{Programmable Logic}
    \acro{PS}{Processing System}
    \acro{RTL}{Register Transfer Level}
    \acro{SoC}{System on Chip}
    \acro{SOM}{System-on-Module}
    \acro{VHDL}{VHSIC Hardware Description Language}
    \acro{VHSIC}{Very High Speed Integrated Circuit}
\end{acronym}
%\end{abbrList}
%\ac{FPGA} případně \acp{FPGA} vygeneruje FGPAs

\chapter{Úvod}

Motivace, definice problému, cíle práce.

\chapter{Nástroje}

\section{AMD Vitis™ HLS for Intuitive Design and Productivity}

\textbf{AMD Vitis™ HLS} je pokročilý vývojový nástroj určený pro automatizovanou syntézu digitálního hardwaru z vysokoúrovňových programovacích jazyků.\\Tento nástroj umožňuje transformaci algoritmického popisu v jazycích C/C++ do nízkoúrovňového \ac{RTL} popisu, typicky ve \ac{VHDL} nebo Verilogu, optimalizovaného pro \ac{FPGA} a adaptivní \ac{SoC} platformy společnosti AMD.
\\
\\
V kontextu návrhu systémů pro zpracování obrazu představuje Vitis HLS klíčový prvek, který abstrahuje komplexitu hardwarového návrhu a umožňuje soustředit se na funkční stránku algoritmů. Nástroj provádí automatickou analýzu kódu, plánování (scheduling) a mapování operací na dostupné hardwarové zdroje (DSP bloky, paměťové bloky).

Mezi klíčové vlastnosti, využívané v této práci, patří:

\begin{itemize}
    \item \textbf{Hardwarová akcelerace:} Generování vysoce optimalizovaných \ac{IP} jader (Intellectual Property cores) s možností definovat výkonnostní parametry, jako je propustnost a latence.
    
    \item \textbf{Optimalizační direktivy:} Využití tzv. pragmas pro řízení architektury výsledného hardwaru, včetně paralelizace smyček (loop unrolling), pipeliningu instrukcí a správy paměťových rozhraní.
    
    \item \textbf{Verifikace:} Integrované prostředí pro C-simulaci (ověření funkčnosti algoritmu) a C/RTL ko-simulaci (ověření bitové shody vygenerovaného \ac{RTL} s původním C++ modelem).

    \item \textbf{Integrace:} Přímá návaznost na sadu AMD Vivado™ Design Suite pro finální implementaci, syntézu a nahrání bitstreamu do cílového zařízení, např. kitu Kria KV260. a  \cite{amd_vitis}
\end{itemize}

\newpage

\section{AMD Vivado™ Design Suite}
\textbf{AMD Vivado™} je komplexní integrované vývojové prostředí pro \ac{SoC} a \ac{FPGA}. Které obsahuje: návrh, syntézu, plánování a mapování, verifikační a simulační nástroje. 

Pro účely této práce slouží Vivado jako centrální integrační platforma, která propojuje hardwarové akcelerátory vyvinuté v nástroji Vitis HLS s procesorovým subsystémem (PS) platformy Kria KV260. Klíčovou komponentou prostředí je nástroj IP Integrator, který umožňuje grafický návrh systému na úrovni bloků (Block Design). Zde jsou jednotlivá IP jádra (Intellectual Property cores) propojena pomocí standardizované sběrnice AMBA® AXI, což zajišťuje efektivní komunikaci mezi programovatelnou logikou a řídicím procesorem ARM.

Proces návrhu ve Vivado Design Suite zahrnuje následující klíčové fáze:

\begin{itemize}
    \item \textbf{RTL Syntéza:} Převod zdrojových kódů (VHDL, Verilog) a blokových schémat do seznamu spojů (netlist).
    
    \item \textbf{Implementace:} Fyzický návrh čipu zahrnující rozmístění (Placement) logických elementů a propojení (Routing) signálových tras. Moderní verze Vivado ML (Machine Learning) využívají v této fázi algoritmy strojového učení pro optimalizaci časování a minimalizaci spotřeby.
    
    \item \textbf{Verifikace:} Ověření funkčnosti návrhu pomocí behaviorální a post-implementační simulace.
    
    \item \textbf{Generování bitstreamu:} Vytvoření konfiguračního souboru pro FPGA část a export hardwarové definice (XSA soubor) pro následný vývoj softwaru.
\end{itemize}

Součástí prostředí je také Vivado Hardware Manager, který umožňuje programování cílového zařízení přes rozhraní JTAG a ladění běžícího hardwaru v reálném čase pomocí integrovaných logických analyzátorů (ILA - Integrated Logic Analyzer).

\section{Vývojová deska kria KV260 Vision AI Starter Kit}

\textbf{xck26-sfvc784-2lv-c
}
AMD Kria™ KV260 Vision AI Starter Kit je navržená jako výuková platforma pro návrh a vývoj pokročilých aplikací počítačového vidění se zaměřením na akceleraci vizuálních algoritmů a strojového učení. Tento kit je založen na architektuře \textbf{Zynq® UltraScale+™ MPSoC} a skládá se z v neprodukčního modulu Kria K26 SOM (System-on-Module) osazeného na nosné desce (carrier card) s bohatou sadou rozhraní pro připojení více kamer a zobrazovacích zařízení pomocí konektoru ON Semiconductor Imager Access System (IAS).
Platforma je určena primárně pro vývoj řešení v oblastech, jako jsou chytrá města, průmyslové vidění (machine vision) a bezpečnostní kamery. Díky integrovanému video kodeku (VCU) a programovatelné logice na FPGA umožňuje hardwarovou akceleraci celého aplikačního řetězce – od snímání obrazu přes jeho zpracování až po inferenci neuronových sítí, a to při zachování nízké spotřeby energie a latence. \cite{amd_kv260}

\begin{table}[htbp]
    \centering
    \caption{Technická specifikace vývojové desky AMD Kria KV260 Vision AI Starter Kit \cite{amd_kv260}}
    \label{tab:kv260-specs}
    \begin{tabular}{@{}ll@{}}
        \toprule
        \textbf{Komponenta / Rozhraní} & \textbf{Specifikace} \\ \midrule
        SoC (Čip) & Zynq™ UltraScale+™ MPSoC (modul K26 SOM) \\
        Obrazový procesor & OnSemi AP1302 13MP ISP\\
        Programovatelná logika & 256 tisíc logických buněk,
144 bloků RAM
64 bloků UltraRAM
1 176 DSP slices\\
        Paměť DRR& 4 GB (4 x 512 Mb x 16 bit)\\
        Kamerová rozhraní & 3$\times$ MIPI CSI-2 (2$\times$ IAS, 1$\times$ RPi) \\
        Video výstupy & 1$\times$ HDMI 1.4, 1$\times$ DisplayPort 1.2a\\
        Síťové rozhraní & 1$\times$ 1Gb Ethernet \\
        Rozšiřující rozhraní & 1$\times$ PMOD (12-pin) \\
        USB & 4$\times$ USB 3.0/2.0 \\
        Úložiště dat & Primární  - 512 Mb QSPI (Queued SPI)
Sekundární - slot pro SDHC kartu\\
    \end{tabular}
\end{table}


\chapter{HW platforma a metodika návrhu}
\section{Platforma Zynq UltraScale+ a Kria KV26}
\subsection{LUT}
(Look-Up Table): Kombinační logika.
\subsection{FF}
(Flip-Flop / Registry): Paměťové prvky, pipelining.
\subsection{LUTRAM}
(Distributed RAM): Distribuovaná paměť pro malé buffery.
\subsection{DSP}
(Digital Signal Processor): Dedikované slice's (řezy/blokyú pro rychlou aritmetiku (MACC operace pro konvoluci a maticové výpočty).
\subsection{BRAM}
Paměťové bloky (Line Buffery, Histogram).



\chapter{Reprezentace obrazových dat a barevné prostory}
\section{Fyzikální podstata obrazu}

Proces zpracování obrazových dat úzce souvisí s fyzikální podstatou světla a způsobem, jakým je vnímáno lidským zrakovým ústrojím a následně technicky reprezentováno v digitálních systémech.

\subsection{Viditelné světlo a jeho vnímání}
Světlo je elektromagnetické záření, přičemž lidské oko je citlivé pouze na jeho velmi úzkou část, označovanou jako viditelné spektrum. Toto spektrum zahrnuje vlnové délky přibližně v rozsahu od 380\,nm do 750\,nm. Různé vlnové délky v tomto intervalu vyvolávají v mozku vjem různých barev – od fialové na spodní hranici (krátké vlnové délky), přes modrou, zelenou, žlutou až po červenou na horní hranici (dlouhé vlnové délky), jak ilustruje obrázek~\ref{fig:cie1931_exact}.

\subsection{Standardní kolorimetrický pozorovatel CIE 1931}
Základem objektivního popisu barev v digitálním zpracování obrazu je model definovaný Mezinárodní komisí pro osvětlování (CIE) v roce 1931. Tento model popisuje průměrné barevné vnímání lidského oka pomocí tří spektrálních funkcí míchání barev: $\bar{x}(\lambda)$, $\bar{y}(\lambda)$ a $\bar{z}(\lambda)$, které odpovídají imaginárním primárním složkám $X, Y, Z$.

Pro libovolné spektrální složení světla $P(\lambda)$ lze vypočítat odpovídající trichromatické složky $X, Y, Z$ integrací (v praxi sumací přes diskrétní kroky, typicky po 5\,nm):

\begin{equation}
X = k \int_{\lambda} P(\lambda) \bar{x}(\lambda) d\lambda, \quad
Y = k \int_{\lambda} P(\lambda) \bar{y}(\lambda) d\lambda, \quad
Z = k \int_{\lambda} P(\lambda) \bar{z}(\lambda) d\lambda
\end{equation}

Kde $Y$ přímo reprezentuje jasovou složku (luminanci) a $k$ je normalizační konstanta. Často se využívá normalizovaných trichromatických souřadnic $(x, y, z)$, které definují pouze chromatičnost (odstín a sytost) nezávisle na jasu, přičemž platí $x+y+z=1$:

\begin{equation}
x = \frac{X}{X+Y+Z}, \quad y = \frac{Y}{X+Y+Z}
\end{equation}

Následná konverze z absolutního prostoru $XYZ$ do konkrétního lineárního RGB prostoru (např. sRGB před gama korekcí) se provádí pomocí standardizované transformační matice $\mathbf{M}$:

\begin{equation}
\begin{bmatrix} R \\ G \\ B \end{bmatrix}_{linear} = \mathbf{M}_{XYZ\to RGB} \begin{bmatrix} X \\ Y \\ Z \end{bmatrix}
\end{equation}

Průběh spektrálních funkcí míchání barev CIE 1931 je znázorněn na obrázku \ref{fig:cie1931_final}.

% --- Vlastní LaTeX/TikZ kód obrázku ---
\begin{figure}[htbp]
    \centering
    \begin{tikzpicture}[xscale=1.4, yscale=1.8]
        % Přepočet X: x = (lambda - 300) / 50
        
        % --- Barevné spektrum (zkráceno přesně na 370nm [x=1.4] až 790nm [x=9.8]) ---
        \shade[left color=black, right color=violet!80]       (1.4, -0.2) rectangle (2.0, -0.05);
        \shade[left color=violet!80, right color=blue]        (2.0, -0.2) rectangle (2.9, -0.05);
        \shade[left color=blue, right color=cyan]            (2.9, -0.2) rectangle (4.0, -0.05);
        \shade[left color=cyan, right color=green]           (4.0, -0.2) rectangle (5.1, -0.05);
        \shade[left color=green, right color=yellow]         (5.1, -0.2) rectangle (5.8, -0.05);
        \shade[left color=yellow, right color=orange]        (5.8, -0.2) rectangle (6.5, -0.05);
        \shade[left color=orange, right color=red]           (6.5, -0.2) rectangle (8.0, -0.05);
        \shade[left color=red, right color=black]            (8.0, -0.2) rectangle (9.8, -0.05);

        % --- Hlavní Osa X (vlnová délka) z 370 nm (1.4) do 800 nm (10.1 pro šipku) ---
        \draw[thick, ->] (1.4, 0) -- (10.1, 0) node[right] {$\lambda$ [nm]};
        
        % --- Levá Osa Y ukotvená na 370 nm (x=1.4) ---
        \draw[thick, ->] (1.4, 0) -- (1.4, 2.2);
        \node[rotate=90] at (0.7, 1.1) {Tristimulové hodnoty CIE 1931};

        % Ticks pro levou osu Y (0 až 2.0)
        \foreach \y/\val in {0/0, 0.5/0.5, 1.0/1.0, 1.5/1.5, 2.0/2.0} {
            \draw (1.35, \y) -- (1.4, \y) node[left, font=\scriptsize] {\val};
        }
        % Pomocná linka y=1.0 
        \draw[dashed, gray!30] (1.4, 1.0) -- (9.8, 1.0);

        % Pravá uzavírací linka na 790 nm (x=9.8)
        \draw[gray!30] (9.8, 0) -- (9.8, 2.1);

        % --- CIE 1931 Color Matching Functions (přesná data dle CIE 15:2018) ---
        
        % Z (Modrá) - Jeden dominantní vrchol, žádný sekundární pík.
        \draw[blue, thick, fill=blue, fill opacity=0.1] 
            plot[smooth, tension=0.6] coordinates {
                (1.6, 0) (2.0, 0.06) (2.4, 0.47) (2.8, 1.75) 
                (3.2, 1.67) (3.6, 0.35) (4.0, 0.27) (4.4, 0.07) (4.8, 0)
            };
        \node[blue, font=\small] at (2.9, 1.95) {$\bar{z}(\lambda)$};

        % Y (Zelená) - Reprezentuje jasovou citlivost oka.
        \draw[green!60!black, thick, fill=green, fill opacity=0.1] 
            plot[smooth, tension=0.6] coordinates {
                (1.6, 0) (3.0, 0.04) (3.6, 0.14) (4.0, 0.32) 
                (4.6, 0.86) (5.1, 1.0) (5.6, 0.87) (6.0, 0.63) 
                (6.6, 0.26) (7.2, 0.06) (8.0, 0.01) (9.6, 0)
            };
        \node[green!60!black, font=\small] at (5.1, 1.15) {$\bar{y}(\lambda)$};

        % X (Červená) - POZOR: Tato křivka má dva vrcholy! 
        % Hlavní vrchol u 595 nm a sekundární vrchol u 440 nm.
        \draw[red, thick, fill=red, fill opacity=0.1] 
            plot[smooth, tension=0.5] coordinates {
                (1.6, 0) (2.0, 0.01) (2.4, 0.13) (2.8, 0.35) 
                (3.2, 0.29) (3.6, 0.09) (4.0, 0.01) (4.4, 0.06) 
                (5.0, 0.43) (5.6, 0.91) (5.9, 1.06) (6.4, 0.85) 
                (7.2, 0.16) (8.0, 0.01) (9.6, 0)
            };
        \node[red, font=\small] at (5.9, 1.22) {$\bar{x}(\lambda)$};

        % --- Značky na ose X (vlastní uspořádání) ---
        \foreach \x/\val in {1.6/380, 2.0/400, 4.0/500, 6.0/600, 8.0/700, 9.6/780} {
            \draw (\x, 0) -- (\x, -0.05) node[below, font=\scriptsize, yshift=-4pt] {\val};
        }

        % --- Hranice viditelnosti (380 a 780 nm) ---
        \draw[dashed, gray!50] (1.6, 0) -- (1.6, 1.8);
        \draw[dashed, gray!50] (9.6, 0) -- (9.6, 1.8);

    \end{tikzpicture}
    \caption{Spektrální funkce míchání barev standardního pozorovatele CIE 1931 (data dle \cite{cie15_2018}). Všimněte si vědecky přesného průběhu funkce $\bar{x}(\lambda)$ (červená), která vykazuje sekundární vrchol v modré oblasti spektra, zatímco funkce $\bar{z}(\lambda)$ (modrá) má vrchol pouze jeden.}
    \label{fig:cie1931_final}
\end{figure}

\begin{figure}[htbp]
    \centering
    \begin{tikzpicture}[xscale=1.4, yscale=1.8]
        % Přepočet X: x = (lambda - 300) / 50
        % Viditelná oblast je 380 nm (x=1.6) až 780 nm (x=9.6).
        % Graf nyní začíná na 370 nm (x=1.4) a končí na 790 nm (x=9.8).
        
        % --- Barevné spektrum (podkladový pruh 380-780) ---
        \shade[left color=black, right color=violet]         (1.4, -0.2) rectangle (2.0, -0.05);
        \shade[left color=violet, right color=blue]          (2.0, -0.2) rectangle (2.9, -0.05);
        \shade[left color=blue, right color=cyan]            (2.9, -0.2) rectangle (4.0, -0.05);
        \shade[left color=cyan, right color=green]           (4.0, -0.2) rectangle (5.1, -0.05);
        \shade[left color=green, right color=yellow]         (5.1, -0.2) rectangle (5.8, -0.05);
        \shade[left color=yellow, right color=orange]        (5.8, -0.2) rectangle (6.5, -0.05);
        \shade[left color=orange, right color=red]           (6.5, -0.2) rectangle (8.0, -0.05);
        \shade[left color=red, right color=black]            (8.0, -0.2) rectangle (9.8, -0.05);

        % --- Hlavní Osa X (vlnová délka) z 370 nm do 800 nm (pro přesah šipky) ---
        \draw[thick, ->] (1.4, 0) -- (10.1, 0) node[right] {$\lambda$ [nm]};
        
        % --- Levá Osa Y ukotvená přesně na 370 nm (x=1.4) ---
        \draw[thick, ->] (1.4, 0) -- (1.4, 2.2);
        \node[rotate=90] at (0.7, 1.1) {Tristimulové hodnoty CIE 1931};

        % Ticks pro levou osu Y (0 až 2.0)
        \foreach \y/\val in {0/0, 0.5/0.5, 1.0/1.0, 1.5/1.5, 2.0/2.0} {
            \draw (1.35, \y) -- (1.4, \y) node[left, font=\scriptsize] {\val};
        }
        % Pomocná linka y=1.0 ukotvená uvnitř nového boxu
        \draw[dashed, gray!30] (1.4, 1.0) -- (9.8, 1.0);

        % Pravá uzavírací linka na 790 nm (x=9.8)
        \draw[gray!30] (9.8, 0) -- (9.8, 2.1);

        % --- CIE 1931 Color Matching Functions (přesné tabulkové uzly) ---
        
        % Z (Modrá) 
        \draw[blue, thick, fill=blue, fill opacity=0.1] 
            plot[smooth, tension=0.6] coordinates {
                (1.6, 0) (2.0, 0.06) (2.4, 0.47) (2.8, 1.75) 
                (3.2, 1.67) (3.6, 0.35) (4.0, 0.27) (4.4, 0.07) (4.8, 0)
            };
        \node[blue, font=\small] at (2.9, 1.95) {$\bar{z}(\lambda)$};

        % Y (Zelená) 
        \draw[green!60!black, thick, fill=green, fill opacity=0.1] 
            plot[smooth, tension=0.6] coordinates {
                (1.6, 0) (3.0, 0.04) (3.6, 0.14) (4.0, 0.32) 
                (4.6, 0.86) (5.1, 1.0) (5.6, 0.87) (6.0, 0.63) 
                (6.6, 0.26) (7.2, 0.06) (8.0, 0.01) (9.6, 0)
            };
        \node[green!60!black, font=\small] at (5.1, 1.15) {$\bar{y}(\lambda)$};

        % X (Červená)
        \draw[red, thick, fill=red, fill opacity=0.1] 
            plot[smooth, tension=0.5] coordinates {
                (1.6, 0) (2.0, 0.01) (2.4, 0.13) (2.8, 0.35) 
                (3.2, 0.29) (3.6, 0.09) (4.0, 0.01) (4.4, 0.06) 
                (5.0, 0.43) (5.6, 0.91) (5.9, 1.06) (6.4, 0.85) 
                (7.2, 0.16) (8.0, 0.01) (9.6, 0)
            };
        \node[red, font=\small] at (5.9, 1.22) {$\bar{x}(\lambda)$};

        % --- Značky na ose X (vlastní uspořádání) ---
        \foreach \x/\val in {1.6/380, 2.0/400, 4.0/500, 6.0/600, 8.0/700, 9.6/780} {
            \draw (\x, 0) -- (\x, -0.05) node[below, font=\scriptsize, yshift=-4pt] {\val};
        }

        % --- Hranice viditelnosti (380 a 780 nm) ---
        \draw[dashed, gray!50] (1.6, 0) -- (1.6, 1.8);
        \draw[dashed, gray!50] (9.6, 0) -- (9.6, 1.8);

    \end{tikzpicture}
    \caption{Pravé hodnoty trichromatických činitelů standardního pozorovatele CIE 1931 v rozmezí vlnových délek 380--780\,nm. Křivka $\bar{y}(\lambda)$ současně reprezentuje poměrnou světelnou účinnost.}
    \label{fig:cie1931_exact}
\end{figure}

Lidské vidění je založeno na trichromatické teorii. Sítnice oka obsahuje světlocitlivé buňky (čípky) tří typů, které jsou citlivé na krátké (S -- modrá), střední (M -- zelená) a dlouhé (L -- červená) vlnové délky. Z těchto biologických předpokladů přímo vychází technické zpracování obrazu.

\subsection{Akvizice barevného obrazu}
Běžné obrazové senzory v digitálních kamerách (CMOS nebo CCD) dokážou měřit pouze intenzitu dopadajícího světla, tj. počet fotonů. Samy o sobě tedy generují pouze monochromatický (černobílý) obraz a nerozlišují vlnovou délku. 

Aby bylo možné zachytit barevný obraz, umísťuje se před jednotlivé pixely senzoru pole barevných filtrů (Color Filter Array -- CFA). Nejrozšířenějším typem je tzv. \textbf{Bayerova maska}, která propouští ke každému pixelu pouze jednu ze tří základních barev (červenou, zelenou, nebo modrou). Zelených filtrů je na senzoru dvojnásobné množství oproti červeným a modrým, což reflektuje vyšší citlivost lidského oka na zelenou barvu. Ze surových dat senzoru se následně plnobarevný obraz s hodnotami RGB pro každý pixel dopočítává pomocí interpolačních algoritmů, přičemž tento proces se nazývá \textit{demosaicing}.

\subsection{Barevný model RGB}
Barevný prostor RGB je aditivní barevný model založený na mísení tří základních (primárních) barev: červené (Red), zelené (Green) a modré (Blue). Tento model je absolutním základem digitálního zobrazování a počítačového vidění, neboť přímo odpovídá fyzikální podstatě činnosti obrazovek a výše zmíněných obrazových senzorů.

\subsubsection*{Princip aditivního míchání}
Kombinací tří primárních složek v různých intenzitách vzniká široké spektrum barev. Při nulové intenzitě všech složek nedochází k žádnému vyzařování, čímž získáváme černou barvu. Naopak při jejich maximální a vyvážené intenzitě vzniká barva bílá. Vlnové délky korespondující se základními složkami modelu jsou definovány přibližně následovně:
\begin{itemize}
    \item \textbf{R (Red):} Odpovídá dlouhým vlnovým délkám ($\sim 630$\,nm).
    \item \textbf{G (Green):} Odpovídá středním vlnovým délkám ($\sim 530$\,nm).
    \item \textbf{B (Blue):} Odpovídá krátkým vlnovým délkám ($\sim 450$\,nm).
\end{itemize}

\subsubsection*{Technická reprezentace}
V počítačové grafice je barva konkrétního obrazového bodu (pixelu) definována uspořádanou trojicí hodnot $(R, G, B)$. Nejčastěji se využívá 8bitová datová hloubka na jeden kanál. Každá složka tedy může nabývat celočíselných hodnot v intervalu $\langle 0, 255 \rangle$. Pro uložení jednoho pixelu v modelu RGB je tak zapotřebí 24 bitů dat (celkem $2^{24} \approx 16,7$ milionu možných barev). V hardwarových implementacích na hradlových polích (FPGA) je tato bitová šířka stěžejní pro návrh propustnosti paměťových sběrnic (např. AXI4-Stream). Přehled vlastností modelu RGB je shrnut v tabulce~\ref{tab:rgb_vlastnosti}.

\begin{table}[htbp]
    \centering
    \caption{Základní charakteristika barevného modelu RGB}
    \label{tab:rgb_vlastnosti}
    \begin{tabular}{@{}ll@{}}
        \toprule
        \textbf{Vlastnost} & \textbf{Popis} \\ \midrule
        Typ modelu & Aditivní (světelný) \\
        Rozsah hodnot & Celočíselně $0\dots255$ (8-bit) nebo normalizovaně $0,0\dots1,0$ \\
        Primární využití & Monitory, televizory, digitální kamery, senzory \\
        Hlavní nevýhoda & Vysoká korelace mezi kanály, neseparovaná složka jasu \\ \bottomrule
    \end{tabular}
\end{table}

\subsubsection*{Standardy a gamuty}
Samotný model RGB představuje pouze abstraktní matematický rámec. Pro praktické využití a zajištění barevné konzistence napříč různými zařízeními musí být definován konkrétní barevný prostor (gamut), který přesně určuje chromatické souřadnice primárních barev a referenční bílý bod. Mezi nejužívanější standardy patří:
\begin{itemize}
    \item \textbf{sRGB:} V současnosti nejrozšířenější standard pro webové aplikace a běžná spotřebitelská zařízení.
    \item \textbf{Adobe RGB:} Nabízí výrazně širší barevný gamut, zejména v oblasti zelených a azurových odstínů, což je výhodné pro profesionální zpracování fotografií a předtiskovou přípravu.
    \item \textbf{ProPhoto RGB:} Velmi rozsáhlý barevný prostor používaný primárně při archivaci a špičkovém postprodukčním zpracování obrazu.
\end{itemize}

\subsubsection*{Omezení modelu RGB při zpracování obrazu}
Ačkoliv je prostor RGB ideální pro zobrazení a akvizici dat, pro složitější počítačové zpracování obrazu nebo kompresi (např. formát JPEG) nebývá optimální. Na rozdíl od prostorů odvozených z modelu YUV (či YCbCr), kde je obraz matematicky rozdělen na jasovou (luminance) a barvovou (chrominance) složku, je v RGB jasová informace úzce provázána a rozprostřena mezi všechny tři kanály. Změna hodnoty jediného kanálu tak ovlivní nejen odstín, ale i celkový jas pixelu. Absence oddělené jasové složky činí formát RGB značně nevhodným pro efektivní datovou kompresi, u které se s výhodou využívá nedokonalosti lidského zraku a jeho nižší citlivosti na změny barev oproti změnám jasu.


\subsection{RGB na HSV}

\[
\begin{aligned}
R' &= \frac{R}{255}, \quad G' = \frac{G}{255}, \quad B' = \frac{B}{255} \\
C_{max} &= \max(R', G', B') \\
C_{min} &= \min(R', G', B') \\
\Delta &= C_{max} - C_{min} \\
\\
H &= \begin{cases} 
0^\circ & \Delta = 0 \\
60^\circ \times \left(\frac{G' - B'}{\Delta} \pmod 6\right) & C_{max} = R' \\
60^\circ \times \left(\frac{B' - R'}{\Delta} + 2\right) & C_{max} = G' \\
60^\circ \times \left(\frac{R' - G'}{\Delta} + 4\right) & C_{max} = B'
\end{cases} \\
\\
S &= \begin{cases} 
0 & C_{max} = 0 \\
\frac{\Delta}{C_{max}} & C_{max} \neq 0 
\end{cases} \\
\\
V &= C_{max}
\end{aligned}
\]



a zpětná transformace 

\[\begin{aligned}
C &= V \times S \\
X &= C \times (1 - |(H / 60^\circ) \pmod 2 - 1|) \\
m &= V - C \\
\\
(R', G', B') &= \begin{cases} 
(C, X, 0) & 0^\circ \le H < 60^\circ \\
(X, C, 0) & 60^\circ \le H < 120^\circ \\
(0, C, X) & 120^\circ \le H < 180^\circ \\
(0, X, C) & 180^\circ \le H < 240^\circ \\
(X, 0, C) & 240^\circ \le H < 300^\circ \\
(C, 0, X) & 300^\circ \le H < 360^\circ
\end{cases} \\
\\
R &= (R' + m) \times 255 \\
G &= (G' + m) \times 255 \\
B &= (B' + m) \times 255 
\end{aligned}\]

\newpage
\section{Barevný prostor YUV a jeho využití}

Lidský zrakový systém (HVS -- \textit{Human Visual System}) je evolučně uzpůsoben tak, že vykazuje podstatně vyšší citlivost na změny jasu (luminance) než na změny barvy (chrominance). Zatímco barevný prostor RGB je ideální pro zobrazení na koncových zařízeních (monitory, displeje), pro účely zpracování, přenosu a komprese obrazu je krajně neefektivní kvůli vysoké korelaci mezi jednotlivými barevnými kanály. Z tohoto důvodu byl definován barevný prostor YUV, který jasovou a barevnou informaci odděluje.

Obecný výpočetní vzorec:
\begin{equation}
\begin{aligned}
    Y = Y_R \cdot R + Y_G \cdot G + Y_B \cdot B \\
    \text{přičemž platí:} \quad Y_R + Y_G + Y_B &= 1 \\[10pt]
    \text{Odvození } P_B \text{ a } P_R \text{:} \qquad P_B &= 1 - Y_B \\
    P_R &= 1 - Y_R \\[10pt]
    C_B &= \frac{B - Y}{2P_B} \\
    C_R &= \frac{R - Y}{2P_R}
\end{aligned}
\end{equation}

\subsection{Definice složek Y, U a V}

Prostor YUV transformuje obraz do jedné jasové a dvou barvonosných (rozdílových) složek. Toto rozdělení vychází z historických potřeb zpětné kompatibility barevného televizního vysílání s černobílými přijímači a je zakotveno ve standardech Mezinárodní telekomunikační unie (ITU).

\begin{itemize}
    \item \textbf{Y (Jasová složka / Luma):} Reprezentuje monochromatickou (černobílou) verzi obrazu. Udává úroveň jasu každého pixelu nezávisle na jeho barvě.
    \item \textbf{U (Chrominance B-Y):} Barevná rozdílová složka reprezentující odchylku modré barvy od celkového jasu.
    \item \textbf{V (Chrominance R-Y):} Barevná rozdílová složka reprezentující odchylku červené barvy od celkového jasu.
\end{itemize}

Rozdílová složka zelené barvy (G-Y) se nepřenáší, jelikož ji lze matematicky odvodit ze složek Y, U a V, čímž se šetří datový tok.

\subsection{Matematické matice převodu z RGB}

Převod mezi prostory RGB a YUV není fixní, nýbrž závisí na váhových koeficientech definujících příspěvek jednotlivých základních barev do celkového jasu. Tyto koeficienty se vyvíjely v závislosti na použitých luminoforech a moderních zobrazovacích technologiích. Ve všech případech se však vychází z normovaných hodnot $R, G, B \in [0, 1]$.

\subsubsection{Standard ITU-R BT.601 (SDTV)}
Tento standard byl primárně definován pro televizi se standardním rozlišením (SDTV). Matice pro dopřednou transformaci z RGB do YUV je definována následovně:

\begin{equation}
\begin{pmatrix}
Y \\
U \\
V
\end{pmatrix}
=
\begin{pmatrix}
0.299 & 0.587 & 0.114 \\
-0.14713 & -0.28886 & 0.436 \\
0.615 & -0.51499 & -0.10001
\end{pmatrix}
\begin{pmatrix}
R \\
G \\
B
\end{pmatrix}
\label{eq:bt601}
\end{equation}

\subsubsection{Standard ITU-R BT.709 (HDTV)}
S nástupem televize s vysokým rozlišením (HDTV) byly koeficienty upraveny tak, aby lépe odpovídaly moderním displejům. Zejména se snížil váhový příspěvek červené a modré složky ve prospěch zelené:

\begin{equation}
\begin{pmatrix}
Y \\
U \\
V
\end{pmatrix}
=
\begin{pmatrix}
0.2126 & 0.7152 & 0.0722 \\
-0.09991 & -0.33609 & 0.436 \\
0.615 & -0.55861 & -0.05639
\end{pmatrix}
\begin{pmatrix}
R \\
G \\
B
\end{pmatrix}
\label{eq:bt709}
\end{equation}

\subsection{Standard ITU-R BT.2020 (UHDTV)}

\begin{equation}
\begin{pmatrix}
Y \\
U \\
V
\end{pmatrix}
=
\begin{pmatrix}
 0.2627 &  0.6780 &  0.0593 \\
-0.1396 & -0.3604 &  0.5000 \\
 0.5000 & -0.4598 & -0.0402
\end{pmatrix}
\begin{pmatrix}
R \\
G \\
B
\end{pmatrix}
\label{eq:bt2020}
\end{equation}

\begin{table}[htbp]
\centering
\caption{Porovnání primárních převodových koeficientů jasové složky pro standardy ITU-R.}
\begin{tabular}{lccc}
\toprule
\textbf{Standard} & $\mathbf{W_R}$ (Váha červené) & $\mathbf{W_G}$ (Váha zelené) & $\mathbf{W_B}$ (Váha modré) \\
\midrule
ITU-R BT.601 (SD)  & 0.2990  & 0.5870  & 0.1140 \\
ITU-R BT.709 (HD)  & 0.2126  & 0.7152  & 0.0722 \\
ITU-R BT.2020 (4K) & 0.2627  & 0.6780  & 0.0593 \\
\bottomrule
\end{tabular}
\label{tab:yuv_coeffs}
\end{table}

\subsection{Diagram U-V barevné roviny}

Vztah mezi jasem a barvou lze vizualizovat v tzv. U-V rovině (při konstantním jasu Y). Na obrázku \ref{fig:uv_plane} je zobrazen vektorový prostor chrominančních složek s vyznačením polohy primárních (RGB) a sekundárních (CMY) barev dle standardu BT.601.

\begin{figure}[htbp]
\centering
\begin{tikzpicture}[scale=4.5, every node/.style={scale=0.9}]
    % Osy
    \draw[->, thick, gray!80] (-0.7, 0) -- (0.7, 0) node[right, text=black] {$U$ (B-Y)};
    \draw[->, thick, gray!80] (0, -0.7) -- (0, 0.7) node[above, text=black] {$V$ (R-Y)};
    
    % Mrizka
    \draw[dashed, gray!30, step=0.2] (-0.6,-0.6) grid (0.6,0.6);
    
    % Body barev BT.601 (Y,U,V)
    % Red: U=-0.147, V=0.615
    \filldraw[red] (-0.147, 0.615) circle (0.6pt) node[above left, text=black] {Červená (R)};
    % Green: U=-0.289, V=-0.515
    \filldraw[green!60!black] (-0.289, -0.515) circle (0.6pt) node[below left, text=black] {Zelená (G)};
    % Blue: U=0.436, V=-0.100
    \filldraw[blue] (0.436, -0.100) circle (0.6pt) node[below right, text=black] {Modrá (B)};
    
    % Cyan: U=0.147, V=-0.615
    \filldraw[cyan] (0.147, -0.615) circle (0.6pt) node[below right, text=black] {Azurová (C)};
    % Magenta: U=0.289, V=0.515
    \filldraw[magenta] (0.289, 0.515) circle (0.6pt) node[above right, text=black] {Purpurová (M)};
    % Yellow: U=-0.436, V=0.100
    \filldraw[yellow!80!black] (-0.436, 0.100) circle (0.6pt) node[above left, text=black] {Žlutá (Y)};
    
    % Pocatek (šedá/bílá/černá - žádná barevná informace)
    \filldraw[black] (0, 0) circle (0.6pt) node[below right] {Neutrální osa (Y)};
    
    % Polygon hranice
    \draw[thick, gray!50] (-0.147, 0.615) -- (0.289, 0.515) -- (0.436, -0.100) -- (0.147, -0.615) -- (-0.289, -0.515) -- (-0.436, 0.100) -- cycle;

\end{tikzpicture}
\caption{Projekce barevného prostoru do 2D chrominanční roviny U-V (tzv. YUV gamut). Neutrální barvy (odstíny šedi) leží v počátku souřadnic (U=0, V=0).}
\label{fig:uv_plane}
\end{figure}

\subsection{Chrominanční podvzorkování (Chroma Subsampling)}

Jedním z nejzásadnějších důvodů pro přechod z RGB do YUV v oblasti zpracování obrazu je možnost provedení tzv. \textit{chrominančního podvzorkování} (chroma subsampling). Vzhledem ke zmíněné nižší citlivosti lidského oka na barvu lze prostorové rozlišení složek U a V redukovat, aniž by pozorovatel zaznamenal zjevnou degradaci obrazové kvality.

Systém zápisu podvzorkování se nejčastěji vyjadřuje poměrem $J:a:b$, kde $J$ je šířka referenčního bloku (obvykle 4 pixely), $a$ je počet vzorků chrominance v prvním řádku a $b$ počet vzorků chrominance ve druhém řádku.

\begin{figure}[htbp]
\centering
\begin{tikzpicture}[scale=0.9, font=\sffamily]

    % Format 4:4:4
    \begin{scope}[xshift=0cm]
        \node[above, font=\bfseries] at (2, 4.5) {4:4:4 (Bez redukce)};
        \foreach \y in {0,1,2,3} {
            \foreach \x in {0,1,2,3} {
                % Y sample
                \draw[thick] (\x, \y) rectangle ++(1,1);
                % U/V sample matches Y
                \fill[blue!20] (\x+0.1, \y+0.1) rectangle ++(0.8,0.8);
                \node at (\x+0.5, \y+0.5) {YUV};
            }
        }
    \end{scope}

    % Format 4:2:2
    \begin{scope}[xshift=5.5cm]
        \node[above, font=\bfseries] at (2, 4.5) {4:2:2 (Redukce 1/2)};
        \foreach \y in {0,1,2,3} {
            \foreach \x in {0,1,2,3} {
                \draw[thick] (\x, \y) rectangle ++(1,1);
                \node at (\x+0.5, \y+0.5) {Y};
            }
            \foreach \x in {0,2} {
                \fill[blue!20, opacity=0.7] (\x, \y) rectangle ++(2,1);
                \node[blue!80!black] at (\x+1, \y+0.5) {U/V};
            }
        }
    \end{scope}

    % Format 4:2:0
    \begin{scope}[xshift=11cm]
        \node[above, font=\bfseries] at (2, 4.5) {4:2:0 (Redukce 3/4)};
        \foreach \y in {0,1,2,3} {
            \foreach \x in {0,1,2,3} {
                \draw[thick] (\x, \y) rectangle ++(1,1);
                \node at (\x+0.5, \y+0.5) {Y};
            }
        }
        \foreach \y in {0,2} {
            \foreach \x in {0,2} {
                \fill[blue!20, opacity=0.7] (\x, \y) rectangle ++(2,2);
                \node[blue!80!black] at (\x+1, \y+1) {U/V};
            }
        }
    \end{scope}

\end{tikzpicture}
\caption{Vizualizace principu chrominančního podvzorkování. Každý čtverec v mřížce představuje jeden pixel. Jasová složka (Y) je zachována vždy v plném rozlišení, zatímco barevná informace (U/V) je sdílena pro více sousedících pixelů.}
\label{fig:subsampling}
\end{figure}

\subsection{Výhody pro kompresi a zpracování obrazu}

Transformace do prostoru YUV poskytuje pro moderní systémy počítačového vidění a kompresní algoritmy (např. JPEG, H.264/AVC, HEVC) následující klíčové benefity:

\begin{enumerate}
    \item \textbf{Datová úspora bez vizuální ztráty:} Aplikací podvzorkování 4:2:0 dochází k redukci celkového objemu dat nekomprimovaného obrazu na polovinu oproti původnímu RGB formátu, a to ještě před aplikací samotných kompresních algoritmů (např. diskrétní kosinové transformace DCT).
    \item \textbf{Energetická efektivita algoritmů:} Mnoho algoritmů zpracování obrazu (detekce hran, sledování pohybu, extrakce příznaků) nevyžaduje barevnou informaci. Namísto konverze celého RGB obrazu na černobílý stačí v systému YUV pouze vyextrahovat existující kanál Y a provést výpočty nad ním.
    \item \textbf{Optimalizace kvantizace:} Během kompresního procesu (např. v normě JPEG) jsou používány rozdílné kvantizační matice pro jas a pro barvu. Matice aplikované na chrominanční kanály odstraňují podstatně více vysokofrekvenčních detailů (agresivnější ztrátová komprese) než u kanálu jasového.
\end{enumerate}


dle standardu \textbf{SDTV ITU-R BT.601}
dle standardu \textbf{SDTV ITU-R BT.709} FULL HP
dle standardu \textbf{SDTV ITU-R BT.2020} 4K (možná až 8?)

\[\begin{aligned}
Y &= 0.299R + 0.587G + 0.114B \\
U &= -0.14713R - 0.28886G + 0.436B \\
V &= 0.615R - 0.51499G - 0.10001B
\end{aligned}\]

zpětný převod dostaneme inverzní maticí

\[\begin{aligned}
R &= Y + 1.13983V \\
G &= Y - 0.39465U - 0.58060V \\
B &= Y + 2.03211U
\end{aligned}\]

\subsection{RGB na YUV}



\subsection{}

\section{Histogram}

\href{https://docs.opencv.org/3.4/d4/d1b/tutorial_histogram_equalization.html}{Histogram opencv}


\section{Filtrace}
Filtrace obrazových dat je jedním s fundamentálních bloků počítáčového vidění. Ať už pro vyhlazení hran, odstranění zašumění nebo vylepšení obrazových dat pro další rozpoznávání. Filtrování založená na operaci konvoluce najdeme prakticky všude.

https://opencv.org/image-filtering-using-convolution-in-opencv/

% Requires: \usepackage{amsmath}
\begin{equation}
    F(x,y)=(f * g)(x, y) = \sum_{m=-\infty}^{\infty} \sum_{n=-\infty}^{\infty} f(m, n)\, g(x - m, y - n)
    \label{eq:placeholder_label}
\end{equation}

Ztráta okrajů a zpoždění (Flushing): Pokud chceme režim SAME (výstupní obraz je stejně velký jako vstupní), okno začne generovat platný středový pixel až ve chvíli, kdy do něj nateče celá polovina jádra. Abychom na konci obrazu ze systému "vytlačili" i ty úplně poslední pixely, musí pipelina běžet o něco déle, než je počet pixelů v obraze. Tomuto se říká Flushing.

Identita
\begin{equation}
    \begin{bmatrix}
        0 & 0 & 0 \\
        0 & 1 & 0 \\
        0 & 0 & 0
    \end{bmatrix}
\end{equation}

\subsection{Průměrovací filtr}
Průměrovací filtr
\begin{equation}
    \frac{1}{9} \cdot 
    \begin{bmatrix}
        1 & 1 & 1 \\
        1 & 1 & 1 \\
        1 & 1 & 1
    \end{bmatrix}
\end{equation}

Zostřovací filtr
\begin{equation}
    \begin{bmatrix}
        0 & -1 & 0 \\
        -1 & 5 & -1 \\
        0 & -1 & 0
    \end{bmatrix}
\end{equation}

\subsection{Hranové filtry}

Prewitt horizontální
\begin{equation}
    \begin{bmatrix}
        -1 & 0 & 1 \\
        -1 & 0 & 1 \\
        -1 & 0 & 1
    \end{bmatrix}
\end{equation}

Prewitt vertikální
\begin{equation}
    \begin{bmatrix}
        -1 & -1 & -1 \\
        0 & 0 & 0 \\
        1 & 1 & 1
    \end{bmatrix}
\end{equation}

Sobel horizontální
\begin{equation}
    \begin{bmatrix}
        -1 & 0 & 1 \\
        -2 & 0 & 2 \\
        -1 & 0 & 1
    \end{bmatrix}
\end{equation}

Sobel vertikální
\begin{equation}
    \begin{bmatrix}
        -1 & -2 & -1 \\
        0 & 0 & 0 \\
        1 & 2 & 1
    \end{bmatrix}
\end{equation}

\subsection{Morfologické operace}

V morfologii se tyto matice označují jako strukturní prvky ($S$).

Čtvercový element ($3 \times 3$)
\begin{equation}
    S_{square} = \begin{bmatrix}
        1 & 1 & 1 \\
        1 & 1 & 1 \\
        1 & 1 & 1
    \end{bmatrix}
\end{equation}

Křížový element ($3 \times 3$)
\begin{equation}
    S_{cross} = \begin{bmatrix}
        0 & 1 & 0 \\
        1 & 1 & 1 \\
        0 & 1 & 0
    \end{bmatrix}
\end{equation}

Aproximace disku ($5 \times 5$)
\begin{equation}
    S_{disk} = \begin{bmatrix}
        0 & 0 & 1 & 0 & 0 \\
        0 & 1 & 1 & 1 & 0 \\
        1 & 1 & 1 & 1 & 1 \\
        0 & 1 & 1 & 1 & 0 \\
        0 & 0 & 1 & 0 & 0
    \end{bmatrix}
\end{equation}


\chapter{Hardwarová architektura a implementace algoritmů}

\section{Datový tok}
AXI video stream
v textu zmínit, že deska podporuje tzv. \textit{overlay} architekturu, což  umožní měnit hardwarovou konfiguraci FPGA bez nutnosti restartu celého systému, to je testování algoritmů zpracování obrazu velmi výhodné. 

\section{Moduly pro barevné transformace}
\subsection{Univerzální barevný převodník pro lineární barevné palety?}

\subsection{Nelineární barevné převodníky}
\subsubsection{RGB na HSV}
\subsection{HSV na RGB}

\section{Metodologie designu}

\subsection{Problém s 2023 a 2025}

Zásadní problém při implementaci ve verzi 2023.2 je nekompletních konfiguračních souborech (Board Files) pro modul Kria KV260, kterým chybělo detailní fyzické mapování hodinových domén (tzv. Clock Regions) na samotném křemíku Zynq UltraScale\bigoplus. Kvůli tomu interní směrovací algoritmus nedokázal správně alokovat trasy pro časově kritické signály, což vedlo ke kritickému selhání NTD (New Timing Driven) flow hned v úvodu fáze Place & Route. Upgrade na verzi Vivado 2025.2 tento problém vyřešil architektonicky, jelikož nová verze nativně integruje aktualizované a plně ověřené profily desky přímo z repozitáře AMD. Tyto nové profily již obsahují bezchybně definované mapování pinů i hodinového směrování, díky čemuž nástroj získal přesnou informaci o fyzické topologii čipu a celá implementace, včetně vygenerování bitstreamu, proběhla úspěšně bez nutnosti definovat dodatečná složitá manuální omezení (constraints).

Problém ve verzi 2025.2 je změna flow z nějakého souboru s zatím nevím jakou přípomou se přešlo na xml, jenžě to zatím není oddláděné a nefuguje jak má.



\textbf{Programovací model Vitis HLS}

Zdrojový kód jazyka C pro Vitis HLS je koncipován tak, aby efektivně využíval výhody a specifické vlastnosti architektury adaptivních SoC a hradlových polí (FPGA) společnosti AMD.

Nástroj Vitis HLS podporuje konstrukty paralelního programování, které umožňují modelovat požadovanou hardwarovou implementaci. Mezi tyto konstrukty patří:
\begin{itemize}
    \item \textbf{HLS tasks (úlohy):} umožňují souběžnost na úrovni procesů (process-level concurrency).
    \item \textbf{HLS vectors (vektory):} umožňují paralelismus na úrovni dat (data-level parallelism).
    \item \textbf{HLS streams (proudy):} umožňují komunikaci mezi souběžně běžícími úlohami.
\end{itemize}
K řízení výsledků syntézy lze využít tzv. \textbf{pragmata} (syntetizační direktivy). Tato pragmata zahrnují instrukce pro zřetězení (pipeline), rozvinutí smyček (unroll), dělení paměťových polí (array partitioning) a definici protokolů rozhraní.

Další podrobnosti naleznete v sekci „HLS Programmers Guide“ v uživatelské příručce \textit{Vitis High-Level Synthesis User Guide}.

 
\begin{figure}
    \centering
    \includegraphics[width=0.5\linewidth]{HLS:methodology.png}
    \caption{Enter Caption}
    \label{fig:placeholder}
\end{figure}
\cite{amd_vitis}

\section{Aritmetika s pevnou řádovou čárkou}
formát Q4.4 Texas insturment nebo ARM
\href{https://nl.mathworks.com/help/simulink/ug/fixed-point-numbers.html}{Fixed point článek - Matlab }

\begin{equation}
    \textbf{Pro formát} QN.N \\
    \text{Největší číslo} {2^n} \\
     \text{Nejmenší číslo} \frac{1}{2^n}
\end{equation}

Příklad převodu
46 * 1 = 
46 * 0,5 = 
46 * 4,625 = 


\begin{table}
    \centering
    \begin{tabular}{|c|c|c|c|c|c|c|c|c|c|l|l|l|l|l|l|l|}\hline
         UQ7.8&  Znaménkový bit&  \multicolumn{7}{|c|}{7-bitů pro celé číslo}&  \multicolumn{8}{|c|}{8-bitů pro desetiné číslo}\\\hline
         Bitová hodnota&  $-2^7$&  $2^6$&  $2^5$&  $2^4$&  $2^3$&  $2^2$&  $2^1$&  $2^0$& $2^{-1}$ & $2^{-2}$& $2^{-3}$& $2^{-4}$& $2^{-5}$& $2^{-6}$&$2^{-7}$ &$2^{-8}$\\ \hline
    \end{tabular}
    \caption{Caption}
    \label{tab:placeholder}
\end{table}

\section{Zpracování video streamu}

Neboli AXI stream video
signály TLAST a TUSER, tedy EOF (End-of-File) a SOF (Start-of-Frame) 
\begin{lstlisting}[language=C]
void video_proc(input_video, output_video){
for(i=0; i < MAX_FRAME_SIZE; i++){
if(VS){
......
}
if(HS){
......
}
if(DE){
//Pixel Processing Algorithm
}
....
}
}
\end{lstlisting}

\section{Převody do barevných soustav}
zminit problém nalezení toho "správného" postupu YCbCr, YUV,  a dalších (
\begin{itemize}
    \item \textbf{Y′UV} is the separation used in PAL.
    \item \href{https://en.wikipedia.org/wiki/YDbDr}{YDbDr} is the format used in \href{https://en.wikipedia.org/wiki/SECAM}{SECAM}, unusually based on non-gamma-corrected (linear) RGB, making the Y component true \href{https://en.wikipedia.org/wiki/Luminance}{luminance}.
    \item \href{https://en.wikipedia.org/wiki/YIQ}{Y′IQ} is the format used in \href{https://en.wikipedia.org/wiki/NTSC}{NTSC} television.
    \item \href{https://en.wikipedia.org/wiki/YPbPr}{Y'PbPr} is the separation used in \href{https://en.wikipedia.org/wiki/Component_video}{component video}.
    \item \href{https://en.wikipedia.org/wiki/YCbCr}{Y′CbCr} is any \href{https://en.wikipedia.org/wiki/Digital_component_video}{digital encoding} of Y'PbPr suited for video and image compression and transmission formats such as \href{https://en.wikipedia.org/wiki/Moving_Picture_Experts_Group}{MPEG} and \href{https://en.wikipedia.org/wiki/JPEG}{JPEG}.\textsuperscript{\href{https://en.wikipedia.org/wiki/Y\%E2\%80\%B2UV\#cite_note-Poynton-2}{[2]}}
    \item \href{https://en.wikipedia.org/wiki/YCoCg}{Y′CoCg} is a similar separation that offers better performance and lossless round-tripping, used in newer video codecs and \href{https://en.wikipedia.org/wiki/JPEG_XR}{JPEG XR}.
\end{itemize}
 ) takže jsem vytvořil univerzální, kde stačí pouze upravit koeficienty

\section{Architektura a hardwarová implementace generického IP jádra pro zpracování obrazu}

\subsection{Úvod do architektury IP jádra}
Navržené IP jádro (Intellectual Property Core) pro konverzi barevných prostorů a bodové operace je koncipováno s důrazem na maximální propustnost, škálovatelnost a efektivní využití hardwarových zdrojů hradlového pole (FPGA). Architektura striktně odděluje fyzické rozhraní modulu (Top-Level) od samotného výpočetního algoritmu. Tento přístup, využívající pokročilých vlastností jazyka C++ (zejména šablonového metaprogramování – *Template Metaprogramming*), umožňuje generovat vysoce optimalizovaný RTL (Register-Transfer Level) kód na míru konkrétní aplikaci bez nutnosti zásahu do zdrojového kódu algoritmu.

\subsection{Rozhraní modulu a datový tok (Dataflow}
Procesorový systém a okolní periferie komunikují s IP jádrem pomocí standardizovaných sběrnic rodiny AMBA AXI4.


\textbf{Datové rozhraní (AXI4-Stream):} Pro kontinuální tok obrazových dat je využito rozhraní AXI4-Stream ($s_axis_video$ pro vstup, $m_axis_video$ pro výstup). Obrazový bod (pixel) je sdružen pomocí direktivy pragma HLS AGGREGATE do jediného datového vodiče (TDATA). Synchronizace snímku je zajištěna postranními signály – začátek snímku (SOF) je indikován signálem TUSER a konec řádku (EOL) signálem TLAST. Díky těmto signálům je jádro navrženo jako datově řízené (\textit{Data-Driven}), což minimalizuje zpoždění a eliminuje nutnost externího stavového automatu (FSM) pro řízení toku.

\textbf{Řídicí rozhraní (AXI4-Lite):} Konfigurační parametry, jakými jsou rozlišení obrazu (width, height), transformační matice (coeffs) a aditivní konstanty (offsets), jsou mapovány do paměťového prostoru procesoru prostřednictvím sběrnice AXI4-Lite. To umožňuje běhové (Run-time) modifikace parametrů algoritmu ze strany operačního systému (např. OS Linux běžící na ARM procesoru).

\subsection{Výpočetní architektura a paralelizace}

Hlavní výpočetní smyčka je navržena pro plné zřetězení instrukcí (\textit{Pipelining}). Aplikací direktivy \#pragma HLS PIPELINE II=1 je syntezátor Vitis HLS donucen vytvořit hardwarovou strukturu, která je schopna přijmout, zpracovat a odeslat jeden pixel v každém hodinovém taktu (Initiation Interval = 1). 

Aby bylo možné dosáhnout propustnosti jednoho pixelu za takt při maticových operacích, bylo nutné překonat omezení čtecích portů standardních blokových pamětí (BRAM). 

\subsubsection{Destrukce paměťových bloků (Array Partitioning)
}
Standardní paměť BRAM v FPGA disponuje pouze dvěma čtecími porty. Maticový výpočet však vyžaduje paralelní přístup ke všem koeficientům (např. 9 koeficientů pro matici 3x3) současně. Z tohoto důvodu byla použita direktiva $\#pragma HLS ARRAY_PARTITION$ variable=coeffs complete. Tato instrukce rozloží logické pole na samostatné klopné obvody (Flip-Flops), čímž vznikne dedikovaný fyzický vodič pro každý koeficient a je zamezeno vzniku úzkého hrdla při čtení dat.

\subsubsection{Využití bloků DSP a ochrana proti přetečení
}
Vnitřní cyklus, iterující přes barevné kanály, je plně rozbalen pomocí direktivy \#pragma HLS UNROLL. Tím je zamezeno vzniku sekvenčního stavového automatu a namísto toho je vytvořen paralelní strom násobiček.


Syntezátor pro tyto operace (Multiply-Accumulate - MAC) automaticky inferuje dedikované DSP řezy (např. DSP48E2 na architektuře UltraScale\+). Pro zajištění numerické stability je akumulační registr definován pomocí datového typu $ap_int<48>$. Šířka 48 bitů přesně mapuje fyzickou architekturu sčítačky uvnitř bloku DSP, čímž je zaručena naprostá odolnost proti přetečení (\textit{Overflow}) i při kaskádním sčítání dílčích součinů. Následná normalizace výpočtu je provedena bitovým posunem vpravo ($\gg 8$), což v hardwaru nepředstavuje žádnou logickou spotřebu, jedná se pouze o přeznačení vodičů.

\subsection{Význam šablonového metaprogramování pro syntézu HW}
Inovativním prvkem navrženého řešení je definice algoritmu ve formě C\+\+ šablony ($template <int IN_CH, int IN_W, int OUT_CH, int OUT_W>$). Tento přístup přináší zásadní výhody pro finální hardwarovou implementaci:

\begin{enumerate}
    \item \textbf{Konstantní propagace v době kompilace:} Parametry šablony (počet kanálů, bitová šířka) jsou vyhodnoceny během kompilace (\textit{Compile-time}). Konstanty pro saturaci obrazu (např. $MAX_VAL$) jsou syntezátorem zadrátovány přímo do komparátorů, čímž odpadá nutnost alokace paměti pro tyto mezní hodnoty.
    \item \textbf{Eliminace mrtvého kódu (Dead-Code Elimination):} Pokud je jádro instanciováno pro konverzi 3kanálového formátu na 1kanálový (např. RGB do Grayscale), HLS nástroj automaticky optimalizuje rozbalené smyčky a fyzicky syntetizuje pouze ty DSP bloky a registry, které jsou pro danou dimenzi matice nezbytné. Neaktivní kanály jsou nástrojem odříznuty, což vede k dramatické úspoře hradlové logiky (LUT) a spotřeby energie na čipu.
    \item \textbf{Univerzalita modulu:} Jeden zdrojový kód pokrývá libovolné konverze (RGB $\rightarrow$ YCbCr, RGBA $\rightarrow$ CMYK, atd.) a podporuje asymetrické bitové šířky vstupu a výstupu, což odpovídá moderním požadavkům metodiky návrhu znovupoužitelných IP jader (IP Reuse).
\end{enumerate}

\subsection{Závěr návrhu}
Kombinací paralelizace na úrovni dat (Data-Level Parallelism), efektivního mapování DSP bloků a generické metaprogramovací architektury bylo vytvořeno robustní IP jádro. Modul bez problémů splňuje požadavky na zpracování video toku v reálném čase (Real-time Video Processing) s minimální latencí, deterministickým časováním a optimální spotřebou zdrojů platformy FPGA.

\section{Passthrouht}
Testovací komponenta, logicky je nejednoduší ale byla důležitá k testování a naučení se ovládání datového streamu, důležitý poznatek je dodržet to co Xillinx uvádí ve své dokumentaci a to je zpracování ve dvou vnořených cyklech, pro výšku a šírku myslím že to je v souboru XAP793 nebo jak se to jmenuje


\section{RGB na HSV}
\subsection{VHDL}

\begin{table}[!h]
    \centering
    \begin{tabular}{cccc}\toprule
         Prostředek&  Syntéza&  Implementace& Procent (\%)\\\midrule
         LUT&  220&  218& 0,19\\
         LUTRAM&  14&  14
& 0,02\\
         FF&  118&  118
& 0,05\\ 
         DSP& 2& 2
&0,16\\
         IO& 58& 58&30,69\\ \bottomrule
    \end{tabular}
    \caption{Srovnání prostředků}
    \label{tab:rgb2hsv}
\end{table}
\subsection{HSL}

\newpage
\section{HSV na RGB}

\subsection{VHDL}

\begin{table}[!h]
    \centering
    \begin{tabular}{cccc}
        Prostředek& Syntéza& Implementace& Procent (\%)\\\midrule
        LUT &  385 & 381 & 0,33 \\
        LUTRAM& 6 & 6 & 0,01 \\
        FF& 181 & 181 & 0,08\\
        DSP& 2 & 2 & 0,16\\
        IO& 58 & 58 & 31,69\\ \bottomrule
    \end{tabular}
    \caption{Srovnání prostředků}
    \label{tab:hsv2rgb}
\end{table}

\subsection{HLS}

\section{2D filter}

\href{https://docs.amd.com/r/en-US/am004-versal-dsp-engine/Filter-Designs}{Filter Designs}

\subsection{Sčítací strom}
\href{https://docs.amd.com/r/en-US/am004-versal-dsp-engine/Adder-Tree}{Adder Tree }

\subsection{Multiplay-Accumulate MACC }
\url{https://docs.amd.com/r/en-US/am004-versal-dsp-engine/MACC-Extension}
Vynásob a sečti

\subsection{HLS}
\subsection{VHDL}

\begin{table}[!h]
    \centering
    \begin{tabular}{cccc}
        Prostředek& Syntéza& Implementace& Procent (\%\(+-0.1\))\\\midrule
        LUT &  11368 & 11355 & 9,71 \\
        LUTRAM& 123 & 123 & 0,21 \\
        FF& 33161 & 33161 & 14,16\\
        DSP& 0 & 0 & 0,0\\
        IO& 98 & 98 & 51,85\\ \bottomrule
    \end{tabular}
    \caption{Srovnání prostředků}
    \label{tab:filter_2d}
\end{table}

 \subsection{Srovnání}

 
\chapter{Závěr}

\printbibliography[title={Použitá literatura},heading=bibintoc]
\end{document}
